\begin{python0}
from solutions import *; clear()
\end{python0}
\sectionthree{Advice}

Do not have too many conversion operators. Doing so might lead to
ambiguities. For instance, suppose you have conversion operators from
\textit{C} to \textit{D} and vice versa; you also have \textit{operator+} in
\textit{C} and \textit{D}. If you run the following code

\begin{consolethree}[escapeinside=||]
C c;
D d;
c + d;
\end{consolethree}

you'll see that your compiler does not know if it should do

\textit{D(c) + d} using \textit{D::operator+}

or

\textit{c + C(d)} using \textit{C::operator+}

Here's another tip: If you have two classes \textit{C} and
\textit{D} where \textit{C} is ``included'' in D, then it's
better to have automatic conversion from \textit{C} to \textit{D}. Consider
the following example:

Let's say you have an \textit{Int} class and a
\textit{Fraction} class. You should have a \textit{operator Fraction()}
declared in the \textit{Int} class. Note that you can achieve the same
thing by having a constructor of the form:

\textit{Fraction(const Int \&) const;}


